% Add to the main preamble (before \begin{document}):
% \usepackage{amsmath,amssymb,booktabs,graphicx,array}
% Input this file only after \maketitle.
% Figures are stored in supplementary/figures/.

\section{Data Processing Details and Dataset Statistics}
\label{sec:supp-data}

\subsection{Data Processing Details}

The main paper provides an overview of the data pipeline;
here, we describe the source-specific processing details.
For Sim7, Real3, and SpatialCog, we extract trajectories
directly from the original videos using a standardized
configuration of 49 RGB frames and 64 query points per
object record.
Tracked pixels are back-projected using depth and camera
intrinsics, and camera poses transform the resulting 3D
points into the first-frame camera coordinate system.
The object mask in the first frame defines the visible
partial surface from which the query points are selected.
RGB frames and point trajectories are generated at the
same temporal indices and therefore require no subsequent
temporal realignment.

For MolmoMotion-1M, we reuse the released trajectories and
associated geometry rather than extracting trajectories
from the source videos.
We transform the trajectories into the same first-frame
camera coordinate system and standardize their point
counts and temporal lengths.
Records with fewer than eight valid object-point candidates
are discarded.
When more than 64 candidates are available, we use
deterministic voxel-uniform sampling to maintain spatial
coverage; for records with 8--63 candidates, we sample
64 points with replacement using a fixed random seed.
RGB frames and point trajectories are aligned using
identical temporal indices.
Intervals longer than 49 frames are uniformly subsampled,
whereas intervals containing 8--48 frames are extended
to 49 frames by repeating the final RGB frame and point
positions.
The resulting padded transitions have zero ground-truth
displacement.
The repeated RGB frames and their zero-displacement
targets are included in training supervision.

\paragraph{Quality control.}
In addition to format standardization, we assess trajectory
quality along four dimensions: geometric validity,
temporal continuity, observation completeness, and motion
plausibility.
Geometric checks reject invalid camera parameters,
non-finite 3D coordinates, and inconsistent depth--position
changes.
Temporal checks detect abrupt changes in point position,
depth, or object extent.
Observation checks limit the fractions of trajectories
that leave the image or lack valid measurements.
Motion checks identify implausible displacement, velocity,
and spatial-scale changes.

The criteria depend on the reliability of each source's
geometry and scale.
For motion magnitude, SpatialCog and Real3 retain peak
displacements in $[0.05,3]$\,m, while metric MolmoMotion-1M
subsets impose an upper limit of 3\,m.
To detect temporal discontinuities, Real3 limits the
99.5th percentile of frame-rate-normalized step displacement
to 0.5\,m; metric MolmoMotion-1M subsets additionally cap
individual steps at 1\,m and speed at 15\,m/s.
DROID also rejects concurrent discontinuities in depth
and 3D position.
For spatial consistency, Real3 and metric MolmoMotion-1M
subsets limit the increase in RMS spatial radius to 0.25\,m.
For observation completeness, these sources cap the
out-of-frame fraction at 5\%; metric MolmoMotion-1M subsets
also cap the unobserved fraction at 5\%.

For YT-VIS, which lacks a reliable metric scale, we instead
use relative step-to-range, adjacent-depth, and spatial-radius
statistics without absolute displacement or speed thresholds.
Simulation records are screened for complete geometry and
valid task annotations.
All motion and continuity checks are evaluated before
temporal padding.

Figure~\ref{fig:supp-data-composition} summarizes the
source-specific candidate counts and retention after
formatting, quality control, and manifest validation,
before the training/evaluation split.
For SpatialCog, 82,831 post-prefilter candidate records
yield 82,588 valid manifest records, corresponding to
a retention rate of 99.7\%.
This pool includes all 400 SpatialCog evaluation samples
in WorldFlowBench.
These samples are excluded at training time, leaving
82,188 SpatialCog records for training.

\begin{figure}[t]
\centering
\includegraphics[width=\linewidth]
{supplementary/figures/data_composition.pdf}
\caption{
Initial data volume and retention across sources before
the training/evaluation split.
(a) Number of candidate object records entering the
reported processing stage, shown on a logarithmic scale.
(b) Fraction retained after formatting, quality control,
and manifest validation, before benchmark exclusion.
For SpatialCog, 82,831 candidate records yield 82,588
valid records, corresponding to a retention rate of 99.7\%.
Initial counts describe source-specific candidate pools
rather than the full original datasets.
Colors indicate the four collections defined in the main
paper.
All counts refer to object records rather than unique videos.
}
\label{fig:supp-data-composition}
\end{figure}

\subsection{Dataset Statistics}

Table~\ref{tab:supp-training-data} summarizes the processed
data inventory. Real3 and SpatialCog exclude WorldFlowBench
records; Sim7 uses its QC inventory before runtime exclusion.
The reported 1,438,379 standardized object clips comprise
70,480,571 frame instances, or approximately 1,223.6 hours
at 16 FPS. With 64 query points and 48 displacement steps
per clip, they contain 92,056,256 point trajectories and
4,418,700,288 3D displacement vectors.
These totals precede Sim7 runtime exclusion and may include
padded frames, repeated query points, and multiple object
records sharing the same source video.

\begin{table}[t]
\centering
\footnotesize
\setlength{\tabcolsep}{3.5pt}
\renewcommand{\arraystretch}{1.12}
\begin{tabular}{@{}ll>{\raggedright\arraybackslash}p{0.36\linewidth}rr@{}}
\toprule
Source
& Collection
& Description
& \shortstack{Object\\clips}
& \shortstack{Standardized\\duration (h)} \\
\midrule
Sim7
& Sim7
& Simulated robot manipulation drawn from seven simulation datasets.
& 190,267 & 161.9 \\
\cmidrule(lr){1-5}

AgiBot World
& Real3
& Real-world robot manipulation across diverse environments and tasks.
& 631,873 & 537.5 \\
\cmidrule(lr){1-5}

STH-V2
& Real3
& Human--object interactions covering a broad range of object motions.
& 86,811 & 73.8 \\
\cmidrule(lr){1-5}

EPIC-KITCHENS
& Real3
& Egocentric human activities and object interactions in kitchens.
& 5,337 & 4.5 \\
\cmidrule(lr){1-5}

SpatialCog
& SpatialCog
& Tabletop and household interactions emphasizing spatial relations.
& 82,188 & 69.9 \\
\cmidrule(lr){1-5}

MolmoSpaces
& MolmoMotion-1M
& Simulated interactive scenes with released 3D motion annotations.
& 164,970 & 140.3 \\
\cmidrule(lr){1-5}

DROID
& MolmoMotion-1M
& Real-world robot manipulation collected across multiple environments
and tasks.
& 19,739 & 16.8 \\
\cmidrule(lr){1-5}

EgoDex
& MolmoMotion-1M
& Egocentric dexterous human--object interactions.
& 126,850 & 107.9 \\
\cmidrule(lr){1-5}

Stereo4D
& MolmoMotion-1M
& Dynamic real-world scenes reconstructed from stereo observations.
& 15,879 & 13.5 \\
\cmidrule(lr){1-5}

Xperience
& MolmoMotion-1M
& Egocentric videos of human activities and object interactions.
& 112,992 & 96.1 \\
\cmidrule(lr){1-5}

YT-VIS
& MolmoMotion-1M
& Diverse Internet videos with non-metric object trajectories.
& 1,473 & 1.3 \\
\midrule
Total
& ---
& ---
& 1,438,379
& 1,223.6 \\
\bottomrule
\end{tabular}
\caption{
Composition of the processed data inventory.
Real3 and SpatialCog counts exclude WorldFlowBench evaluation
records; Sim7 uses its QC inventory before runtime benchmark
exclusion.
Each object clip contains 49 frames standardized to 16 FPS.
Standardized duration in hours is computed as
$N_{\mathrm{clips}} \times 49 / (16 \times 3600)$,
where $N_{\mathrm{clips}}$ denotes the reported number of
object clips.
This measures standardized data volume rather than
deduplicated source-video duration, since different object
records may share the same underlying video interval.
}
\label{tab:supp-training-data}
\end{table}

The MolmoMotion-1M sources exhibit substantial variation
in their native temporal coverage.
Figure~\ref{fig:supp-native-length} reports their interval
lengths before conversion to the common 49-frame
representation and the fraction of records requiring
tail padding.

\begin{figure}[t]
\centering
\includegraphics[width=\linewidth]
{supplementary/figures/native_trajectory_lengths.pdf}
\caption{
Native temporal coverage of the six retained MolmoMotion-1M
subsets before conversion to 49 frames.
(a) Native interval-length distributions; thin lines
indicate P5--P95, thick bars indicate the interquartile
range, and circles indicate the median.
The dashed line marks the target length of 49 frames.
(b) Percentage of records requiring tail padding.
Intervals longer than 49 frames are uniformly subsampled.
}
\label{fig:supp-native-length}
\end{figure}

To characterize the scale of the motion targets, we define
the component RMS of record $r$ as
\begin{equation}
m_r
=
\left(
\frac{1}{48\cdot64\cdot3}
\sum_{t=1}^{48}
\sum_{i=1}^{64}
\left\lVert\mathbf{f}_{r,i,t}\right\rVert_2^2
\right)^{1/2}.
\label{eq:supp-motion-rms}
\end{equation}
Figure~\ref{fig:supp-motion-statistics} reports the
distribution of $m_r$ within each source together with
the source-level aggregate RMS,
$\sqrt{\mathbb{E}_r[m_r^2]}$.
These statistics are computed before flow normalization
and include zero-motion transitions introduced by temporal
padding.
They characterize the numerical scale of the motion targets
rather than physical speed or endpoint displacement.

\begin{figure}[t]
\centering
\includegraphics[width=\linewidth]
{supplementary/figures/motion_statistics.pdf}
\caption{
Pre-normalization motion-scale statistics.
Left: median, 90th percentile, and 99th percentile of
the per-record component RMS.
Right: source-level aggregate component RMS.
Both horizontal axes use logarithmic scales.
SpatialCog statistics use 82,188 records after excluding
all 400 evaluation samples; Real3 statistics also exclude
benchmark records.
Sim7, marked with an asterisk, uses its QC inventory before
runtime benchmark exclusion and provides aggregate moments
but no per-record percentile summary.
Statistics include zero-motion transitions introduced by
temporal padding.
YT-VIS values follow the non-metric scale of its annotations.
}
\label{fig:supp-motion-statistics}
\end{figure}